\section{Introduction}

FunSearch~\cite{funsearch}, Evolution of Heuristics (EoH)~\cite{eoh},
and ReEvo~\cite{reevo} combine language-model program generation with
performance-guided search. Their success raises a narrower question:
do parent conditioning and numerical performance feedback improve the
final heuristics beyond independent generation? Improvements during
search can reflect adaptation to the training instances rather than
better decisions on unseen problems.

I study this question on the minimum weighted dominating set (MWDS)
problem, where weighted greedy provides an interpretable reference and
exact optimization certifies quality on small graphs. COVER-EVO holds
construction, pruning and feasibility checking fixed while the LLM
proposes vertex-scoring expressions. Independent generation (IND),
parent-only evolution (EVO-P), and evolution with numerical prompt
feedback (EVO-F) are compared with parameter search and handcrafted
controls. Each branch develops its own archive and parent history,
so the comparisons concern complete proposal policies.

The evaluation addresses three questions.
\textbf{RQ1:} Do the evolutionary proposal policies improve held-out quality
over independent generation?
\textbf{RQ2:} How do the frozen heuristics transfer across graph distributions,
and what do feature interventions reveal about their behavior?
\textbf{RQ3:} How do solution quality and execution efficiency compare with
non-LLM controls?
The contribution is a controlled comparison of proposal policies,
an evaluation of frozen-program transfer, and an analysis of how
scoring changes survive construction and pruning. Synthetic shifts
separate cost allocation from topology, while a published MWDS
benchmark tests external transfer.

\section{Background and MWDS}

\subsection{Weighted Domination as Set Cover}

Let $G=(V,E)$ be an undirected graph with $n=|V|$ vertices
and positive integer costs $c_v$. The closed neighborhood is
$N[v]=\{v\}\cup\{u:(u,v)\in E\}$.
MWDS seeks a subset $D\subseteq V$ minimizing
$C(D)=\sum_{v\in D}c_v$ such that
$\bigcup_{v\in D}N[v]=V$.
An equivalent integer formulation is
\begin{align}
\min_{\boldsymbol{x}}\quad
    &\sum_{v\in V}c_vx_v, \label{eq:ilp}\\
\text{s.t.}\quad
    &\sum_{v\in N[u]}x_v\geq 1,
    &&\forall u\in V, \notag\\
    &x_v\in\{0,1\},
    &&\forall v\in V. \notag
\end{align}

Each closed neighborhood is a set with selection cost $c_v$, making MWDS
an instance of weighted set cover. Given uncovered vertices $U$, standard
weighted greedy selects a vertex maximizing $|N[v]\cap U|/c_v$.
It achieves an approximation ratio of $H_d\leq 1+\ln d$ for
$d=\max_{v\in V}|N[v]|$~\cite{chvatal}. Removing redundant selected
vertices while preserving coverage retains this guarantee. Arbitrary
learned scoring functions have no such guarantee.

\subsection{Related Work and Design Rationale}

AlphaEvolve and LLaMEA~\cite{alphaevolve,llamea} allow more flexible
algorithmic generation. Restricting COVER-EVO to scoring expressions
keeps construction and feasibility fixed across proposal strategies,
but excludes changes to the broader solver. This is a design trade-off
within heuristic generation~\cite{burke}.
Classical methods such as FastMWDS~\cite{fastmwds} combine reductions
and local search for large-scale weighted domination; the objective is
total cost rather than cardinality~\cite{weighted}.
Reported generalization failures for LLM-evolved bin-packing
heuristics~\cite{sim2025} motivate testing frozen programs beyond their
training distribution. A handcrafted rarity-weighted coverage priority,
RGR+P, provides an interpretable control within the same scaffold.
Section~\ref{sec:rarity_form} develops its mathematical interpretation.
